% !TEX root = ../main.tex
\begin{ouverture}[infrarouge]{8}{Molécules organiques et spectroscopie infrarouge}{Partie 1 · Constitution et transformations de la matière}
{Formule semi-développée\enspace·\enspace Squelette carboné\enspace·\enspace Groupe caractéristique\enspace·\enspace Nomenclature\enspace·\enspace Infrarouge}
\objectif{Passer d'une formule brute à une formule semi-développée, et inversement.}
\objectif{Identifier les groupes caractéristiques des alcools, aldéhydes, cétones et acides carboxyliques.}
\objectif{Justifier le nom d'une molécule simple à partir de sa formule semi-développée, et inversement.}
\objectif{Exploiter un spectre infrarouge à l'aide d'une table de nombres d'onde.}
\objectif{Visualiser la géométrie de molécules organiques avec des modèles ou un logiciel.}
\prerequis{Schémas de Lewis et géométrie des molécules (chapitre~6).}
\prerequis{Règles du duet et de l'octet ; liaisons simples et doubles (chapitre~6).}
\end{ouverture}
\essentiel
\essSection{1}{Écrire la formule d'une molécule organique}
\begin{definition}[Formules brute et semi-développée]
Une molécule organique contient des atomes de carbone et d'hydrogène, souvent d'oxygène. Sa \textbf{formule brute} indique le
nombre d'atomes de chaque élément : $\ce{C2H6O}$ pour l'éthanol. Sa \textbf{formule semi-développée} montre l'enchaînement des
atomes, sans représenter les liaisons avec les atomes d'hydrogène : $\ce{CH3-CH2-OH}$.
\end{definition}
\begin{definition}[Squelette carboné saturé]
L'enchaînement des atomes de carbone forme le \textbf{squelette carboné}. Il est \textbf{saturé} si tous ses atomes de carbone ne
sont liés que par des liaisons simples. Il peut être \textbf{linéaire}, \textbf{ramifié} ou \textbf{cyclique}. Les alcanes, de formule
$\mathrm{C}_n\mathrm{H}_{2n+2}$, n'ont que des liaisons C--C et C--H : méthane, éthane, propane, butane, pentane, hexane pour $n=1$ à 6.
\end{definition}
\essSection{2}{Groupes caractéristiques et familles}
\begin{propriete}[Quatre familles]
Un \textbf{groupe caractéristique} est un groupe d'atomes qui confère des propriétés communes aux molécules qui le portent :
celles-ci forment une \textbf{famille}. Dans un aldéhyde, l'atome de carbone du groupe carbonyle est lié à un atome
d'hydrogène, au bout de la chaîne ; dans une cétone, il est lié à deux atomes de carbone.
\end{propriete}
\schema{08-familles}
\essSection{3}{Nommer une molécule}
\begin{savoirfaire}[Du nom à la formule, et inversement]
\begin{questions}
\item Repérer la \textbf{chaîne principale} : la plus longue chaîne carbonée qui contient le carbone du groupe caractéristique.
Son nombre d'atomes donne le préfixe : méth- (1), éth- (2), prop- (3), but- (4), pent- (5), hex- (6).
\item Numéroter la chaîne pour donner au groupe caractéristique le plus petit numéro possible.
\item Ajouter la terminaison de la famille : \textbf{-ol} pour un alcool (propan-2-ol), \textbf{-al} pour un aldéhyde (propanal),
\textbf{-one} pour une cétone (butanone), et « acide \dots\textbf{oïque} » pour un acide carboxylique (acide éthanoïque).
\item Nommer les ramifications en préfixe, avec le numéro de l'atome qui les porte. Ainsi, la molécule
\[\ce{CH3-CH(CH3)-CH2-OH}\]
est le 2-méthylpropan-1-ol.
\end{questions}
\end{savoirfaire}
\essSection{4}{La spectroscopie infrarouge}
\begin{definition}[Spectre infrarouge]
Une molécule absorbe certains rayonnements infrarouges, qui font vibrer ses liaisons. Un \textbf{spectre infrarouge} représente
la transmittance $T$, en \%, en fonction du \textbf{nombre d'onde} $\sigma=\dfrac{1}{\lambda}$, en \unit{cm^{-1}}, porté de droite à gauche. Chaque
liaison absorbe dans un domaine caractéristique : une bande vers le bas révèle sa présence.
\end{definition}
\begin{propriete}[Table des bandes caractéristiques]
\par\smallskip\noindent\hfill\begin{tabular}{lll}
\toprule
Liaison & Nombre d'onde (\unit{cm^{-1}}) & Allure de la bande \\
\midrule
O--H d'un alcool, libre (en phase gazeuse) & 3580 à 3670 & fine \\
O--H d'un alcool, lié (liquide) & 3200 à 3400 & forte et large \\
O--H d'un acide carboxylique & 2500 à 3200 & forte et très large \\
C--H & 2800 à 3000 & moyenne à forte \\
C--H d'un aldéhyde & 2650 à 2850 & deux bandes moyennes \\
C=O & 1650 à 1750 & forte et fine \\
\bottomrule
\end{tabular}\hfill\null\par\smallskip
En dessous de \qty{1500}{cm^{-1}}, de nombreuses bandes forment l'« empreinte digitale » de la molécule, difficile à interpréter.
\end{propriete}
\schema{08-spectres-ir}
\begin{python}
table = {'O-H alcool': (3200, 3400), 'O-H acide': (2500, 3200),
         'C-H': (2800, 3000), 'C=O': (1650, 1750)}
pics = [3340, 2970, 1715]          # en cm-1

for pic in pics:
    liaisons = [nom for nom, (mini, maxi) in table.items() if mini <= pic <= maxi]
    print(pic, liaisons)
\end{python}
\begin{console}
3340 ['O-H alcool']
2970 ['O-H acide', 'C-H']
1715 ['C=O']
\end{console}
\noindent\capacite{Capacité numérique} Pour chaque nombre d'onde relevé sur un spectre, le programme liste les liaisons dont le domaine le
contient. À \qty{2970}{cm^{-1}}, deux liaisons sont possibles : seule l'allure de la bande, fine ou très large, permet de trancher.
\begin{attention}
Une bande C=O ne suffit pas à distinguer un aldéhyde d'une cétone : il faut chercher les deux bandes C--H de l'aldéhyde,
vers \qty{2720}{cm^{-1}} et \qty{2820}{cm^{-1}}. Dans un nom, le numéro du groupe caractéristique est omis lorsqu'il n'y a pas d'ambiguïté :
on écrit propanone et butanone, mais pentan-2-one.
\end{attention}
\begin{experience}[Visualiser des molécules organiques]
Avec un logiciel de visualisation moléculaire, on affiche les modèles éclatés de l'éthanol, du propanal, de la propanone
et de l'acide éthanoïque. On constate que les atomes de carbone liés par des liaisons simples sont entourés de quatre
liaisons orientées vers les sommets d'un tétraèdre, tandis que l'atome de carbone d'un groupe C=O est au centre d'un
triangle plan.
\end{experience}
\begin{rediger}
Pour exploiter un spectre, relever les bandes au-dessus de \qty{1500}{cm^{-1}}, avec leur nombre d'onde et leur allure, les attribuer à
l'aide de la table, puis conclure sur la famille. Signaler aussi l'absence d'une bande attendue : c'est souvent un argument
décisif.
\end{rediger}
\exercices

\rubrique{Découverte}
\begin{blocLivre}{ex-8-1}
\begin{exercice}[1]{Formules brutes et semi-développées}
\begin{questions}
\item Écrire la formule semi-développée et la formule brute de l'éthanol, dont la formule développée comporte une chaîne
de deux atomes de carbone et un groupe $\ce{-OH}$ en bout de chaîne.
\item Écrire la formule brute de la propanone $\ce{CH3-CO-CH3}$ et de l'acide méthanoïque $\ce{HCOOH}$.
\item Écrire la formule développée du butane $\ce{CH3-CH2-CH2-CH3}$, puis sa formule brute.
\item Vérifier que la formule brute du butane est de la forme $\mathrm{C}_n\mathrm{H}_{2n+2}$.
\end{questions}
\end{exercice}
\end{blocLivre}
\begin{blocLivre}{ex-8-2}
\begin{exercice}[1]{Reconnaître les familles}
Identifier le groupe caractéristique et la famille de chaque molécule.
\begin{questions}
\item $\ce{CH3-CH2-CHO}$ et $\ce{HCHO}$.
\item $\ce{CH3-CO-CH2-CH3}$.
\item $\ce{CH3-CH(OH)-CH3}$ et $\ce{CH3-CH2-CH2-CH2-OH}$.
\item $\ce{CH3-CH2-COOH}$.
\end{questions}
\end{exercice}
\end{blocLivre}
\begin{blocLivre}{ex-8-3}
\begin{exercice}[1]{Nommer des molécules}
Nommer chaque molécule en justifiant.
\begin{questions}
\item $\ce{CH3-CH2-CH2-OH}$ et $\ce{CH3-CH(CH3)-CH2-OH}$.
\item $\ce{CH3-CH2-CO-CH2-CH3}$.
\item $\ce{CH3-CH(CH3)-CHO}$.
\item $\ce{CH3-CH2-CH2-COOH}$.
\end{questions}
\end{exercice}
\end{blocLivre}
\begin{blocLivre}{ex-8-4}
\begin{exercice}[1]{Du nom à la formule}
Écrire la formule semi-développée et la formule brute de chaque molécule, puis indiquer sa famille.
\begin{questions}
\item Hexan-2-one.
\item Acide propanoïque.
\item Méthanal et butanal.
\item 3-méthylbutan-1-ol.
\end{questions}
\end{exercice}
\end{blocLivre}
\begin{blocLivre}{ex-8-5}
\begin{exercice}[1]{Lire un spectre infrarouge}
Un composé X a pour formule brute $\ce{C3H8O}$. Voici son spectre infrarouge.
\schema{08-ex05}
\begin{questions}
\item Relever les bandes situées au-dessus de \qty{1500}{cm^{-1}} et les attribuer à l'aide de la table de l'essentiel.
\item À quelle famille X appartient-il ? Quelle bande attendue est absente ?
\item Parmi le propan-1-ol, le propanal et la propanone, lesquels ont pour formule brute $\ce{C3H8O}$ ?
\item Proposer une formule semi-développée pour X. Le spectre permet-il de distinguer le propan-1-ol du propan-2-ol ?
\end{questions}
\end{exercice}
\end{blocLivre}
\rubrique{Entraînement}
\begin{blocLivre}{ex-8-6}
\begin{exercice}[2]{Attribuer des bandes en Python}
On reprend le programme de l'essentiel.
\begin{questions}
\item \capacite{Capacité numérique}Prévoir l'affichage pour \code{pics = [3000, 2720, 1730]}.
\item Quel défaut du programme la ligne affichée pour \qty{2720}{cm^{-1}} révèle-t-elle ? Compléter le dictionnaire \code{table}.
\item Prévoir alors la ligne affichée pour \qty{2720}{cm^{-1}}.
\item Quelle information du spectre le programme n'utilise-t-il pas ? Comment pourrait-on l'ajouter ?
\end{questions}
\end{exercice}
\end{blocLivre}
\begin{blocLivre}{ex-8-7}
\begin{exercice}[2]{Deux isomères}
Le propanal et la propanone ont la même formule brute : ce sont des isomères.
\begin{questions}
\item Écrire leurs formules semi-développées et leur formule brute commune.
\item Entourer et nommer leur groupe caractéristique ; indiquer leur famille.
\item Pourquoi n'écrit-on pas « propan-2-one » ni « propan-1-al » ?
\item Quelles bandes infrarouges permettent de les distinguer ?
\end{questions}
\end{exercice}
\end{blocLivre}
\begin{blocLivre}{ex-8-8}
\begin{exercice}[2]{Quatre alcools}
Quatre alcools ont pour formule brute $\ce{C4H10O}$ :
A : $\ce{CH3-CH2-CH2-CH2-OH}$ ; B : $\ce{CH3-CH(OH)-CH2-CH3}$ ; C : $\ce{CH3-CH(CH3)-CH2-OH}$ ; D : $\ce{CH3-C(CH3)(OH)-CH3}$.
\begin{questions}
\item Vérifier qu'ils ont tous la même formule brute.
\item Nommer A et B.
\item Nommer C et D, qui ont un squelette ramifié.
\item Leurs spectres infrarouges présentent tous une large bande vers \qty{3300}{cm^{-1}}. Le spectre permet-il de les distinguer ?
\end{questions}
\end{exercice}
\end{blocLivre}
\begin{blocLivre}{ex-8-9}
\begin{exercice}[2]{Du vin au vinaigre}
Au contact de l'air, des bactéries transforment l'éthanol du vin en acide éthanoïque : le vin « tourne » au vinaigre.
\begin{questions}
\item Écrire les formules semi-développées de l'éthanol et de l'acide éthanoïque. Entourer leurs groupes caractéristiques.
\item À l'aide des spectres de l'essentiel, citer deux différences entre les spectres de ces deux molécules.
\item Comment le spectre d'un vin qui commence à tourner évolue-t-il au fil des jours ?
\item Pourquoi la bande O--H de l'acide est-elle beaucoup plus large que celle de l'alcool ?
\end{questions}
\end{exercice}
\end{blocLivre}
\begin{blocLivre}{ex-8-10}
\begin{exercice}[2]{Des molécules odorantes}
L'hexanal sent l'herbe coupée, l'acide butanoïque le beurre rance, l'heptan-2-one le fromage à pâte persillée et le
3-méthylbutan-1-ol la fermentation.
\begin{questions}
\item Écrire la formule semi-développée de chaque molécule.
\item Indiquer la famille de chacune.
\item Écrire leurs formules brutes.
\item Laquelle de ces molécules présente une bande C=O, sans bande C--H d'aldéhyde ni bande O--H ?
\end{questions}
\end{exercice}
\end{blocLivre}
\begin{blocLivre}{ex-8-11}
\begin{exercice}[2]{Des modèles moléculaires}
\schema{08-ex11}
\begin{questions}
\item Écrire la formule brute de chaque molécule modélisée.
\item Écrire leurs formules semi-développées et les nommer.
\item Quelle est la géométrie autour de l'atome de carbone central du modèle A ? autour du carbone lié à l'oxygène dans B ?
\item Quelle bande infrarouge présente l'un des spectres et pas l'autre, au-dessus de \qty{3000}{cm^{-1}} ?
\end{questions}
\end{exercice}
\end{blocLivre}
\rubrique{Objectif DS}
\begin{blocLivre}{ex-8-12}
\begin{exercice}[3]{Oxyder le propan-1-ol}
\begin{docu}{1}{Une oxydation}
L'oxydation du propan-1-ol par les ions permanganate en milieu acide peut conduire à deux produits : un aldéhyde, si
l'oxydant est en défaut et le produit distillé au fur et à mesure, ou un acide carboxylique, si l'oxydant est en excès.
\end{docu}
\begin{docu}{2}{Les spectres de deux produits obtenus}
\schema{08-ex12}
\end{docu}
\begin{questions}
\item Écrire la formule semi-développée du propan-1-ol et justifier son nom.
\item Écrire les formules semi-développées et les noms de l'aldéhyde et de l'acide carboxylique susceptibles de se former.
\item Attribuer, en justifiant, chaque spectre à l'un des deux produits.
\item \initiative Comment vérifier, à partir du spectre de chaque produit, qu'il ne reste plus de propan-1-ol ?
\end{questions}
\end{exercice}
\end{blocLivre}
\begin{blocLivre}{ex-8-13}
\begin{exercice}[3]{Le gel hydroalcoolique}
\begin{docu}{1}{Deux formules}
L'Organisation mondiale de la santé propose deux formulations de solution hydroalcoolique : l'une contient \qty{80}{\percent} en volume
d'éthanol, l'autre \qty{75}{\percent} en volume de propan-2-ol, complétés avec de l'eau, un peu d'eau oxygénée et de glycérol.
\end{docu}
\begin{docu}{2}{Données}
Masses volumiques : éthanol \qty{0,789}{g.mL^{-1}} ; propan-2-ol \qty{0,786}{g.mL^{-1}}. Masses molaires (\unit{g.mol^{-1}}) : H 1,0 ; C 12,0 ; O 16,0.
\end{docu}
\begin{questions}
\item Écrire les formules semi-développées de l'éthanol et du propan-2-ol, et justifier leurs noms.
\item Le propan-2-ol a un isomère qui est aussi un alcool : le nommer et écrire sa formule.
\item Calculer la quantité d'éthanol dans \qty{100}{mL} de la première solution et celle de propan-2-ol dans \qty{100}{mL} de la seconde.
\item \initiative Le spectre infrarouge d'un gel du commerce présente une bande très large et intense vers \qty{3300}{cm^{-1}}. Ce spectre
permet-il de savoir quel alcool il contient ?
\end{questions}
\end{exercice}
\end{blocLivre}
\begin{blocLivre}{ex-8-14}
\begin{exercice}[3]{Identifier un composé inconnu}
\begin{docu}{1}{Une formule brute}
Une analyse élémentaire montre qu'un liquide incolore, utilisé comme solvant de peintures, a pour formule brute $\ce{C4H8O}$
et que son squelette carboné est saturé. On sait aussi qu'il contient un groupe carbonyle.
\end{docu}
\begin{docu}{2}{Son spectre infrarouge}
\schema{08-ex14}
\end{docu}
\begin{questions}
\item Écrire les formules semi-développées et les noms des trois composés carbonylés saturés de formule $\ce{C4H8O}$.
\item Relever les bandes caractéristiques du spectre.
\item \initiative Identifier le composé, en justifiant par la présence et l'absence de bandes.
\item Dessiner son schéma de Lewis et indiquer la géométrie autour de l'atome de carbone du groupe carbonyle.
\end{questions}
\end{exercice}
\end{blocLivre}
\begin{blocLivre}{ex-8-15}
\begin{exercice}[3]{Suivre une oxydation par infrarouge}
\begin{docu}{1}{La transformation}
On oxyde du propan-2-ol en propanone. En fin d'expérience, on isole la phase organique et l'on enregistre son spectre
infrarouge, que l'on compare à celui du réactif.
\end{docu}
\begin{docu}{2}{Les spectres}
\schema{08-ex15}
\end{docu}
\begin{questions}
\item Écrire les formules semi-développées du réactif et du produit ; indiquer leurs familles.
\item Quelle bande apparaît au cours de la transformation ? Laquelle diminue ?
\item \initiative La transformation du propan-2-ol est-elle complète ? Justifier à partir des spectres.
\item Proposer une méthode pour séparer le produit du réactif restant, sachant que la propanone bout à \qty{56}{\degreeCelsius} et le
propan-2-ol à \qty{82}{\degreeCelsius}.
\end{questions}
\end{exercice}
\end{blocLivre}
